% Diseño del libro impreso (solo PDF). Lo copia scripts/qmd.py al proyecto
% Quarto y Quarto lo inserta en el preámbulo de libro.tex, de modo que el
% .tex sigue compilando solo desde build/quarto/.
%
% La paleta es la de la portada tipográfica (portada.svg del EPUB): tinta
% azul noche, acento verde azulado y su versión clara.

\usepackage{tikz}
\usepackage[most]{tcolorbox}
\usepackage{scrlayer-scrpage}

\definecolor{tinta}{HTML}{071B2B}
\definecolor{acento}{HTML}{0B7480}
\definecolor{acentoclaro}{HTML}{69D3C5}
\definecolor{crema}{HTML}{F4F0DF}
\definecolor{gris}{HTML}{4A5560}
\definecolor{fondosigno}{HTML}{EAF5F3}
\definecolor{alerta}{HTML}{A63D2A}
\definecolor{fondoalerta}{HTML}{FBEFEA}

\newfontfamily\fuenteportada{Libertinus Serif Display}

% --- Numeración: parte → capítulo → ficha. El cuerpo de la ficha no se
% numera: «1.3.4.4 Dónde NO confiar» era ruido, no orientación.
\setcounter{secnumdepth}{\subsectionnumdepth}

% --- Figuras ----------------------------------------------------------------
% Cada ficha abre página con título, enlace, figura y pie. La figura se topa
% en el 62 % de la altura útil para que esos cuatro quepan juntos: con el tope
% de Quarto (la página entera) las figuras verticales saltaban a la página
% siguiente y dejaban el título solo. La misma macro de Quarto, con el tope.
\makeatletter
\renewcommand*\pandocbounded[1]{%
  \sbox\pandoc@box{#1}%
  \Gscale@div\@tempa{0.62\textheight}{\dimexpr\ht\pandoc@box+\dp\pandoc@box\relax}%
  \Gscale@div\@tempb{\linewidth}{\wd\pandoc@box}%
  \ifdim\@tempb\p@<\@tempa\p@\let\@tempa\@tempb\fi
  \ifdim\@tempa\p@<\p@\scalebox{\@tempa}{\usebox\pandoc@box}%
  \else\usebox{\pandoc@box}%
  \fi}
\makeatother

% Página con base libre: con `flushbottom` LaTeX estiraba el espacio entre el
% órgano y su primera ficha y el título bajaba a media página.
\raggedbottom

% Sin líneas viudas ni huérfanas.
\clubpenalty=10000
\widowpenalty=10000
\displaywidowpenalty=10000

% --- Encabezados -----------------------------------------------------------
\KOMAoptions{chapterprefix=false}
\setkomafont{disposition}{\sffamily\color{tinta}}
\setkomafont{chapter}{\sffamily\fontsize{30}{36}\selectfont\color{tinta}}
\setkomafont{section}{\sffamily\large\color{acento}}
\setkomafont{subsection}{\sffamily\LARGE\bfseries\color{tinta}}
\setkomafont{subsubsection}{\sffamily\normalsize\bfseries\color{acento}}
\setkomafont{paragraph}{\sffamily\normalsize\bfseries\color{tinta}}

\renewcommand*{\chapterformat}{%
  {\fontsize{64}{64}\selectfont\color{acento}\thechapter}\par\nobreak
  \vspace{4pt}{\color{acentoclaro}\rule{28mm}{2.5pt}}\par\nobreak\vspace{14pt}}
\renewcommand*{\chapterlinesformat}[3]{#2#3}
% La parte abre en página propia; el filtro salta de página tras ella.
\RedeclareSectionCommand[beforeskip=55mm, afterskip=12mm]{chapter}

% El capítulo temático u órgano es una antesala de sus fichas: rótulo
% espaciado sobre una línea fina.
\renewcommand*{\sectionformat}{\thesection\enskip\textperiodcentered\enskip}
% Órgano y ficha llevan línea; las secciones internas de la ficha, no.
\makeatletter
\renewcommand*{\sectionlinesformat}[4]{%
  \Ifstr{#1}{subsubsection}{\@hangfrom{\hskip #2#3}{#4}}{%
  \Ifstr{#1}{paragraph}{\@hangfrom{\hskip #2#3}{#4}}{%
  \parbox{\linewidth}{\raggedright#3#4\par\vspace{2pt}%
  {\color{acentoclaro}\rule{\linewidth}{0.6pt}}}}}}
\makeatother
\RedeclareSectionCommand[beforeskip=4mm, afterskip=5mm]{section}

\renewcommand*{\subsectionformat}{{\color{acento}\thesubsection}\enskip}
\RedeclareSectionCommand[beforeskip=2mm, afterskip=5mm]{subsection}
\RedeclareSectionCommand[beforeskip=5mm, afterskip=1.5mm]{subsubsection}

% --- Cabeceras y folios ----------------------------------------------------
\pagestyle{scrheadings}
\clearpairofpagestyles
\automark[subsection]{chapter}
\ohead{\headmark}
\ofoot*{\pagemark}
\setkomafont{pageheadfoot}{\sffamily\footnotesize\color{gris}}
\setkomafont{pagenumber}{\sffamily\footnotesize\color{acento}}
\KOMAoptions{headsepline=0.4pt}
\setkomafont{headsepline}{\color{acentoclaro}}

% --- Índice ----------------------------------------------------------------
\addtokomafont{chapterentry}{\sffamily\color{tinta}}
\addtokomafont{chapterentrypagenumber}{\color{acento}}
\DeclareTOCStyleEntry[entryformat=\sffamily\color{acento}]{tocline}{section}
\DeclareTOCStyleEntry[entryformat=\small, pagenumberformat=\small]{tocline}{subsection}

% --- Bloques de la ficha (los abre diseno-pdf.lua) ---------------------------
\newtcolorbox{triadasigno}{enhanced, breakable, colback=fondosigno,
  colframe=acento, boxrule=0pt, leftrule=3pt, arc=0pt, outer arc=0pt,
  left=10pt, right=10pt, top=7pt, bottom=7pt, before skip=8pt,
  after skip=10pt, parbox=false}
\newtcolorbox{alertaficha}{enhanced, breakable, colback=fondoalerta,
  colframe=alerta, boxrule=0pt, leftrule=3pt, arc=0pt, outer arc=0pt,
  left=10pt, right=10pt, top=7pt, bottom=7pt, before skip=4pt,
  after skip=10pt, parbox=false}
\newcommand{\etiquetasigno}[1]{{\sffamily\bfseries\color{acento}#1}}
\newenvironment{evidencia}{\par\small}{\par}
\newenvironment{enlaceficha}{\par\small\sffamily\color{gris}\raggedright}{\par}

% --- Portadilla y créditos --------------------------------------------------
\newenvironment{portadillapdf}{%
  \thispagestyle{empty}\null\vspace*{30mm}%
  \begin{flushleft}
  {\fuenteportada\fontsize{34}{42}\selectfont\color{tinta}\titulolineas\par}
  \vspace{6mm}{\color{acentoclaro}\rule{28mm}{2.5pt}}\par\vspace{8mm}
  \sffamily\color{tinta}}
  {\end{flushleft}\clearpage}
\newenvironment{creditospdf}{\small}{\clearpage}

% --- Portada: la misma composición que portada.svg ------------------------
\newcommand{\portadaatlas}{%
  \thispagestyle{empty}%
  \begin{tikzpicture}[remember picture, overlay,
      x=\paperwidth/1600, y=-\paperwidth/1600,
      shift={(current page.north west)}]
    \fill[tinta] (0,0) rectangle (1600,2260);
    \fill[acento] (0,1620) .. controls (420,1440) and (720,1810) .. (1110,1615)
      .. controls (1320,1510) and (1450,1480) .. (1600,1545)
      -- (1600,2260) -- (0,2260) -- cycle;
    \draw[acentoclaro, line width=20*\paperwidth/1600] (1260,320) circle (150);
    \node[anchor=base west, inner sep=0, text=crema,
          align=left, font=\fuenteportada\fontsize{44}{55}\selectfont]
      at (120,580) {\titulolineas};
    \node[anchor=base west, inner sep=0, text=acentoclaro,
          font=\fuenteportada\fontsize{21}{26}\selectfont] at (125,905)
      {\subtitulolibro};
    \node[anchor=base west, inner sep=0, text=crema,
          font=\sffamily\fontsize{15}{18}\selectfont] at (125,1915)
      {\autorlibro};
    \node[anchor=base west, inner sep=0, text=crema!80!acentoclaro,
          font=\sffamily\fontsize{11}{14}\selectfont] at (125,2000)
      {\editoriallibro\enskip\textperiodcentered\enskip\versionlibro};
  \end{tikzpicture}%
  \clearpage}

\AtBeginDocument{%
  \renewcommand{\maketitle}{\portadaatlas}%
  \captionsetup{font={small,sf}, labelfont={bf,color=acento},
    labelsep=period, justification=raggedright, singlelinecheck=false}%
  \hypersetup{colorlinks=true, linkcolor=tinta, urlcolor=acento,
    citecolor=acento}%
}
